\chapter*{Kolofon}
\addcontentsline{toc}{chapter}{Kolofon}
\markboth{Kolofon}{Kolofon}

Naskah ini ditulis dalam bahasa Indonesia, disusun menjadi enam bagian dan tiga puluh
dua bab, dan dibangun dengan \LaTeX{} (pdf\LaTeX{}, kelas \texttt{book}). Seluruh
sumber naskah dalam bentuk \LaTeX{} disertakan bersama buku ini, sehingga setiap angka
dan setiap potongan kode dapat diperiksa ulang pembaca; versi tiga belas bab yang
mendahuluinya tetap diarsipkan sebagai catatan sejarah penyusunan.

Buku ini karya tulis orisinal: kalimat, urutan penyajian, contoh hitung, dan latihan
disusun sendiri. Gagasan matematisnya bersumber pada kitab-kitab di Daftar Rujukan,
dan pembaca didorong membacanya. Tidak ada bagian naskah yang disalin verbatim dari
kitab sumber; bila sebuah definisi atau teorema mengikuti sumber secara ketat,
sumbernya disebut pada peta bab.

Setiap angka di dalam Lab komputasi berasal dari eksekusi nyata di mesin penyusun
(Python 3, NumPy, SciPy) dan dapat diulang pembaca. Potongan kode Lean~4 di Bab~28
dijalankan dengan Lean 4.10.0 tanpa Mathlib dan keluarannya dikutip apa adanya;
taktik pustaka seperti \texttt{linarith} dan \texttt{ring} tidak dijalankan karena
Mathlib tidak terpasang, dan bab itu menyatakan batas tersebut.

Sampul dan tata letak sederhana dengan sengaja, tanpa foto atau gambar raster: seluruh
gambar di dalam bab digambar sebagai kode TikZ, sehingga pembaca dapat memeriksa
angkanya dan mengubahnya, dan daftar gambarnya ada di halaman depan. Enam kotak
belajar menjadi penanda visual utama --- \emph{Untuk pembaca muda}, \emph{Catatan untuk
dosen}, \emph{Untuk peneliti}, \emph{Bacakan keras}, \emph{Kerjakan dengan tangan},
\emph{Resep langkah}, dan \emph{Tantangan} --- sehingga satu bab yang sama dapat
dimasuki dari cara belajar yang berbeda.

\subsection*{Lisensi}
Teks buku ini dilisensikan Creative Commons BY-NC-SA 4.0: boleh dibagikan dan
diadaptasi untuk kepentingan nonkomersial dengan atribusi dan lisensi yang sama.
Potongan kode Python dan Lean~4 di dalamnya dilisensikan MIT, agar dapat dipakai ulang
tanpa hambatan. Jika kamu memerlukan izin di luar ketentuan itu, hubungi penulis.

\subsection*{Cara mengutip}
Pangimpian, Afiq Andico. \emph{Kenapa Mesin Bisa Salah? Dari Aljabar Linear sampai
Keamanan Machine Learning}. Edisi Pertama. 2026.

\subsection*{Versi dan umpan balik}
Versi naskah 1.0, Oktober 2026. Setiap kekeliruan yang kamu temukan --- salah hitung,
rujukan yang tidak cocok, atau kalimat yang menyesatkan --- lebih berguna daripada
pujian; kirimkan ke \texttt{afiqandico13@gmail.com} dengan nomor bab dan halaman.
